\section{Grundlagen, Modell und Daten}

\subsection{Mechanistische Interpretierbarkeit}

Die meiste Forschung zu Sprachmodellen beschränkt sich auf die Beziehung zwischen Eingabe und Ausgabe: Man beobachtet, was ein Modell tut, ohne zu wissen, wie es dazu kommt. Die \gls{mechanisticinterpretability} geht einen Schritt weiter und öffnet die sprichwörtliche Black Box -- ein Bild, das den Unterschied gut greifbar macht, ist der zwischen dem Fahren eines Autos und dem Verstehen seines Motors. Erst das Verständnis der inneren Abläufe erlaubt es, Fehler vorherzusagen und gezielt einzugreifen \cite{nanda2023comprehensive}.

\subsection{Residual Stream und fünf Werkzeuge}

Moderne \glspl{transformer} bestehen aus einem Stapel von Schichten, die alle auf einen gemeinsamen Informationskanal zugreifen, den \emph{Residual Stream}: Jede Schicht liest daraus, berechnet einen Beitrag und addiert ihn wieder hinzu, ohne den bisherigen Inhalt zu löschen \cite{elhage2021mathematical}. Man kann sich das wie ein Notizblatt vorstellen, auf das jede Schicht etwas ergänzt, statt es zu überschreiben -- am Ende wird aus diesem Notizblatt die Vorhersage abgelesen.

Die Arbeit nutzt darauf aufbauend fünf Werkzeuge, die von beschreibenden hin zu kausalen Aussagen führen. Die \textbf{Residual-Stream-Analyse} und das \textbf{lineare Probing} prüfen, ob Sentiment in den \glspl{activation} vorhanden und linear dekodierbar ist, also ob schon eine einfache Trennlinie reicht, um positiv von negativ zu unterscheiden. Der \textbf{Logit Lens} projiziert Zwischenrepräsentationen einzelner Schichten in den Vokabularraum \cite{nostalgebraist2020logit} und zeigt so, ab wann das Modell eine Information nutzt. Die \textbf{\gls{attention}-Analyse} zeigt, welche Köpfe auf welche Token zugreifen, also wohin Information \emph{fließt}. Nur das \textbf{Activation Patching} weist nach, dass eine Komponente kausal \emph{benötigt} wird, indem einzelne \glspl{activation} gezielt ersetzt und die Wirkung auf die Vorhersage gemessen wird \cite{meng2022locating, goldowskydill2023localizing}. Die ersten vier Werkzeuge liefern Korrelationen, nur das letzte liefert einen echten kausalen Nachweis.

Für die Sentiment-Bewertung wird durchgehend das Opinion-Lexikon von Hu und Liu verwendet, das positive und negative Wörter wie „good“/„bad“ oder „excellent“/„terrible“ enthält \cite{hu2004mining}.

\subsection{Modell: Pythia-410M}

Pythia-410M ist ein offenes, vollständig dokumentiertes Decoder-Only-\gls{transformer}-Modell (GPT-NeoX-Architektur) von EleutherAI mit rund 405 Millionen Parametern \cite{biderman2023pythia}. Diese verteilen sich auf vier Hauptkomponenten: Embedding- und Unembedding-Matrix tragen je rund 12,7~Prozent, die \gls{attention}-Blöcke über alle 24 Schichten rund 24,9~Prozent, und die Feed-Forward-Netzwerke vereinen mit rund 49,7~Prozent fast die Hälfte der Parameter auf sich -- der größte Teil der Modellkapazität liegt also nicht in der Attention selbst, sondern in den Feed-Forward-Blöcken. Tabelle~\ref{tab:arch} fasst die zentralen Architekturwerte zusammen.

\begin{table}[htbp]
  \centering
  \caption{Zentrale Architekturwerte von Pythia-410M.}
  \label{tab:arch}
  \begin{tabular}{ll}
    \toprule
    Eigenschaft & Wert \\
    \midrule
    Architekturtyp & Decoder-Only (GPT-NeoX) \\
    Vokabulargröße & 50.304 Token \\
    Hidden Size & 1.024 \\
    Transformer-Schichten & 24 \\
    \glspl{attentionhead} je Schicht & 16 \\
    Positionskodierung & Rotary Positional Embeddings (RoPE) \\
    \bottomrule
  \end{tabular}
\end{table}

\subsection{Datensätze}

Drei Datensätze ergänzen sich, jeder mit einer eigenen Rolle im roten Faden der Arbeit. Das \textbf{Hu-\&-Liu-Opinion-Lexikon} liefert 2.254 Ein-Token-Wörter (865 positiv, 1.389 negativ), die der Tokenizer von Pythia-410M als genau ein Token darstellt -- nur für solche Wörter lässt sich eindeutig ein Embedding-Vektor zuordnen, weshalb sie die Grundlage für Embedding-Analysen und lineares Probing bilden \cite{hu2004mining}. Das \textbf{Sentiment Challenge Dataset} enthält 636 binär gelabelte, sprachlich gezielt schwierige Sätze mit Annotationen wie Negation, Ironie, Idiomen oder gemischtem Sentiment und dient der verhaltensbasierten Bewertung in Kapitel 3 \cite{barnes2019sentiment}. Die \textbf{Counterfactually Augmented Data (CAD)} stellen 1.707 Paare aus einem positiven und einem nahezu kontextgleichen negativen Text bereit -- weil sich beide Texte nur in wenigen sentimenttragenden Wörtern unterscheiden, eignen sie sich besonders gut, um Sentiment isoliert zu betrachten, und werden für die aggregierten Logit-Lens- und \gls{attention}-Analysen in den Kapiteln 5 und 6 genutzt \cite{kaushik2020learning}.
